% !TEX root = AImath.v4.tex

\chapter{人工智能的数学工具}

\begin{introduction}
  \item 微积分在机器学习中的核心应用
  \item 线性代数与高维数据处理
  \item 概率统计在不确定性建模中的作用
  \item 优化理论与算法设计的数学基础
  \item 信息论与数据压缩的数学原理
\end{introduction}

在前一章中，我们探讨了人工智能发展中的三种核心思维：数学思维、算法思维和工程思维。这三种思维为我们提供了从理论到实践的完整框架。然而，要真正理解和应用这些思维，我们需要掌握支撑它们的具体数学工具。

本章将深入介绍人工智能中最重要的数学工具。这些工具不仅是{\bf 数学思维}的具体体现，也是{\bf 算法思维}转化为可计算方法的基础，更是{\bf 工程思维}实现高效系统的数学保障。从微积分的梯度优化，到线性代数的高维数据处理，从概率统计的不确定性建模，到优化理论的算法设计，这些数学工具构成了现代人工智能的理论基石。

\section{预备知识}

人工智能的数学基础涵盖了多个数学分支，每个分支都为解决特定类型的问题提供了理论框架和计算方法。

\subsection{微积分}
微积分是人工智能中最基础的数学工具之一，特别是在机器学习和深度学习中：
- 导数和偏导数：用于计算梯度，是优化算法的核心
- 链式法则：反向传播算法的数学基础
- 积分：用于概率密度函数的计算和期望值的求解
- 多元微积分：处理高维优化问题

\subsection{概率统计}
概率统计为处理不确定性和数据分析提供了理论基础：
- 概率分布：描述随机变量的行为模式
- 贝叶斯定理：用于概率推理和机器学习中的参数估计
- 统计推断：从样本数据推断总体特征
- 假设检验：验证模型的有效性

\subsection{优化理论}
优化理论是机器学习算法的核心：
- 梯度下降法：最常用的优化算法
- 拉格朗日乘数法：处理约束优化问题
- 凸优化：保证全局最优解的优化问题
- 随机优化：处理大规模数据的优化方法

\subsection{信息论}
信息论为量化信息和不确定性提供了数学框架：
- 熵：衡量信息的不确定性
- 互信息：衡量变量间的相关性
- KL散度：比较概率分布的差异
- 信息增益：决策树算法的核心概念

\subsection{计算理论}
计算理论研究算法的复杂性和可计算性：
- 时间复杂度：算法运行时间的分析
- 空间复杂度：算法内存使用的分析
- NP完全性：困难问题的分类
- 近似算法：处理NP困难问题的方法

\subsection{离散数学}
离散数学为算法设计和分析提供了基础：
- 组合数学：用于分析和解决离散结构的问题，如组合优化、排列和组合
- 布尔代数：用于逻辑推理和数理逻辑，特别是在逻辑编程和电路设计中
- 图论和网络流：用于优化和分析网络结构

\section{线性代数基础}

要学习深度学习，首先需要掌握一些基本技能。所有机器学习方法都涉及从数据中提取信息。因此，我们先学习一些关于数据的实用技能，包括存储、操作和预处理数据。

机器学习通常需要处理大型数据集。我们可以将某些数据集视为一个表，其中表的行对应样本，列对应属性。线性代数为人们提供了一些用来处理表格数据的方法。我们不会太深究细节，而是将重点放在矩阵运算的基本原理及其实现上。

深度学习是关于优化的学习。对于一个带有参数的模型，我们想要找到其中能拟合数据的最好模型。在算法的每个步骤中，决定以何种方式调整参数需要一点微积分知识。本章将简要介绍这些知识。幸运的是，`autograd`包会自动计算微分，本章也将介绍它。

机器学习还涉及如何做出预测：给定观察到的信息，某些未知属性可能的值是多少？要在不确定的情况下进行严格的推断，我们需要借用概率语言。

最后，官方文档提供了本书之外的大量描述和示例。在本章的结尾，我们将展示如何在官方文档中查找所需信息。

本书对读者数学基础无过分要求，只要可以正确理解深度学习所需的数学知识即可。但这并不意味着本书中不涉及数学方面的内容，本章会快速介绍一些基本且常用的数学知识，以便读者能够理解书中的大部分数学内容。如果读者想要深入理解全部数学内容，可以进一步学习本书数学附录中给出的数学基础知识。

\subsection{矩阵的应用现状}

矩阵在机器学习和深度学习上面的应用，有以下几个方面：

1. 在将数据集应用于机器学习模型之前，需要对数据集做预处理，比如给数据去噪、降维，这样可以减少计算资源和提高模型预测准确度。常用的去噪降维的方法是主成分分析（PCA），里面就用到了矩阵，具体步骤为：（1）求解数据集的协方差矩阵的特征值和特征向量，（2）根据特征值从大到小排序（3）选取前面几个特征值对应的特征向量作为新的坐标基向量。（4）将所有的数据集映射到这个新的坐标系中，从而得到降噪降维之后的数据集了。在第（3）步中只选取前面最大的几个特征值所对应的特征向量，是因为小特征值所蕴含的信息量比较小，它们很可能是一些噪声数据。

2. 另一个在机器学习中的应用是奇异值分解，奇异值分解的公式为M=U∑V，在Netflix Prize电影推荐系统中，根据用户对电影的评分来建模的推荐系统使用到了奇异值分解的思想。

3. 在深度学习模型的框架中，数据是以矩阵的形式组织的，比如一次输入64张彩色图片，输入数据是64x3x32x32维度的。这样一方面可以简化表示，另一方面可以提高计算速度，特别是GPU能够针对这个进行优化，如果不以矩阵的形式表示数据，那么GPU的优势就没发挥出来。

矩阵是人工智能基础范畴中不可或缺的一部分。必备的数学知识是理解人工智能不可或缺的要素，所有的人工智能技术归根到底都建立在数学模型之上，而这些数学模型又都离不开线性代数的理论框架。

线性代数不仅仅是人工智能的基础，更是现代数学和以现代数学作为主要分析方法的众多学科的基础。从量子力学到图像处理都离不开向量和矩阵的使用。而在向量和矩阵背后，线性代数的核心意义在于提供了一种看待世界的抽象视角：万事万物都可以被抽象成某些特征的组合，并在由预置规则定义的框架之下以静态和动态的方式加以观察。线性代数是AI专家必须掌握的知识。

如果你只想把人工智能和机器学习的工具当作一个黑匣子，那么你只需要足够的数学计算就可以确定你的问题是否符合模型使用。如果你想提出新想法，线性代数则是你必须要学习的东西。并不是说你需要学习有关数学的所有知识，这样会耽搁于此，失去研究其他更重要的东西（如微积分/统计）的动力。

线性代数是应用数学领域，这是人工智能专家离不开的。如果你不精通这个领域，你永远不会成为一个好的AI专家。正如斯凯勒·斯皮克曼所说："线性代数是21世纪的数学。"线性代数有助于产生新的想法，这就是为什么它是人工智能科学家和研究人员必须学习的东西。他们可以用标量、向量、张量、矩阵、集合和序列、拓扑、博弈论、图论、函数、线性变换、特征值和特征向量的概念来建立模型。

\subsection{矩阵基本运算}

\subsubsection{加减法}
对应位置相加减。例如：
$$\begin{pmatrix}1&0\\0&1\end{pmatrix}+\begin{pmatrix}0&1\\1&0\end{pmatrix}=\begin{pmatrix}1&1\\1&1\end{pmatrix}$$
$$\begin{pmatrix}1&0\\0&1\end{pmatrix}-\begin{pmatrix}0&1\\1&0\end{pmatrix}=\begin{pmatrix}1&-1\\-1&1\end{pmatrix}$$
需要注意的是，这两个矩阵必须行列数均相同。

\subsubsection{矩阵数乘}
矩阵的所有元素都乘上这个数即可。
$$k\cdot\begin{pmatrix}0&1\\1&0\end{pmatrix}=\begin{pmatrix}0&k\\k&0\end{pmatrix}$$

\subsubsection{矩阵乘法}
两个矩阵的乘法仅当第一个矩阵A的列数和另一个矩阵B的行数相等时才能定义。如果矩阵A是m×n矩阵（m行n列），矩阵B是n×k矩阵（n行k列），则它们的乘积C是一个m×k矩阵。对于矩阵C的每个元素，它如下计算：

$$c_{i,j}=a_{i,1} b_{1,j}+a_{i,2} b_{2,j}+\cdots+a_{i,n} b_{n,j}=\sum_{k=1}^n a_{i,k} b_{k,j}$$

容易看出，单位矩阵乘任何矩阵，不论左乘还是右乘，都不会改变矩阵。这在传统的线性代数教材中已经被证明很多次了，这里不再赘述。

\subsubsection{矩阵转置}
将每个元素的行数和列数互换即可。
$$\begin{pmatrix}1&2\\3&4\end{pmatrix}^T=\begin{pmatrix}1&3\\2&4\end{pmatrix}$$

\subsubsection{共轭矩阵}
共轭是复数域的概念，只是把矩阵中每个元素取其共轭的复数即可。
$$A=\begin{pmatrix}3+i&5\\2-2i&i\end{pmatrix}, \quad \overline{A}=\begin{pmatrix}3-i&5\\2+2i&-i\end{pmatrix}$$

\subsubsection{矩阵的特征值及特征向量}
设A是n阶方阵，如果数λ和n维非零列向量x使关系式Ax=λx成立，那么这样的数λ称为矩阵A特征值，非零向量x称为A的对应于特征值λ的特征向量。对于矩阵的特征值及特征向量，在传统线性代数教材中已经有了足够充分的解释，并对其性质和计算方法也有了足够多的介绍，这里不再赘述了。矩阵的特征值和特征向量可以揭示线性变换的深层特性，我们稍后会展开介绍其具体应用，并简述其实现原理。

\subsubsection{线性变换}
在数学上，如果满足下式，那么映射F(a)就是线性的：
F(a + b) = F(a) + F(b) 以及 F(ka) = kF(a)

如果映射F保持了基本运算：加法和数量乘，那么就可以称该映射为线性的。在这种情况下，将两个向量相加然后再进行变换得到的结果和先分别进行变换再将变换后的向量相加得到的结果相同。同样，将一个向量数量乘再进行变换和先进行变换再数量乘的结果也是一样的。

这个线性变换的定义有两条重要的引理：
(1) 映射F(a) = aM，当M为任意方阵时，此映射是一个线性变换，这是因为：
F(a + b)= (a + b)M = aM + bM = F(a)+ F(b)
F(ka) = (ka)M = k(aM) = kF(a)

(2) 零向量的任意线性变换的结果仍然是零向量。（如果F(0) = a，a ≠ 0。那么F不可能是线性变换。因为F(k0) = a，但F(k0) ≠ kF(0)，因此线性变换不会导致平移（原点位置上不会变化）。

\subsubsection{矩阵的仿射变换}
仿射变换是指线性变换后接着平移。因此仿射变换的集合是线性变换的超集，任何线性变换都是仿射变换，但不是所有仿射变换都是线性变换。

任何具有形式$v^\prime = vM + b$的变换都是仿射变换。

\subsubsection{矩阵的可逆变换}
如果存在一个逆变换可以"撤销"原变换，那么该变换是可逆的。换句话说，如果存在逆变换G，使得G(F(a)) = a，对于任意a，映射F(a)是可逆的。

存在非仿射变换的可逆变换，但暂不考虑它们。现在，我们集中精力于检测一个仿射变换是否可逆。一个仿射变换就是一个线性变换加上平移，显然，可以用相反的量"撤销"平移部分，所以问题变为一个线性变换是否可逆。

显然，除了投影以外，其他变换都能"撤销"。当物体被投影时，某一维有用的信息被抛弃了，而这些信息时不可能恢复的。因此，所有基本变换除了投影都是可逆的。因为任意线性变换都能表达为矩阵，所以求逆变换等价于求矩阵的逆。如果矩阵是奇异的，则变换不可逆；可逆矩阵的行列式不为0。

\subsubsection{矩阵的等角变换}
如果变换前后两向量夹角的大小和方向都不改变，该变换是等角的。只有平移，旋转和均匀缩放是等角变换。等角变换将会保持比例不变，镜像并不是等角变换，因为尽管两向量夹角的大小不变，但夹角的方向改变了。所有等角变换都是仿射和可逆的。

\subsubsection{矩阵的正交变换}
术语"正交"用来描述具有某种性质的矩阵。正交变换的基本思想是轴保持互相垂直，而且不进行缩放变换。

平移、旋转和镜像是仅有的正交变换。长度、角度、面积和体积都保持不变。（尽管如此，但因为镜像变换被认为是正交变换，所以一定要密切注意角度、面积和体积的准确定义）。

正交矩阵的行列式为1或者负1，所有正交矩阵都是仿射和可逆的。

\section{图像处理中的数学应用}

图像处理是人工智能中一个重要的应用领域，它将抽象的数学概念转化为具体的视觉处理技术。

\subsection{图像的数字化表示}

\subsubsection{灰度图像}
在电子计算机领域中，灰度（Gray scale）数字图像是每个像素只有一个采样颜色的图像。这类图像通常显示为从最暗黑色到最亮的白色的灰度，尽管理论上这个采样可以是任何颜色的不同深浅，甚至可以是不同亮度上的不同颜色。灰度图像与黑白图像不同，在计算机图像领域中黑白图像只有黑白两种颜色，灰度图像在黑色与白色之间还有许多级的颜色深度。

灰度图像经常是在单个电磁波频谱如可见光内测量每个像素的亮度得到的。用于显示的灰度图像通常用每个采样像素8 bits的非线性尺度来保存，这样可以有256种灰度（8bits就是2的8次方=256）。这种精度刚刚能够避免可见的条带有损，并且非常易于编程。在医学图像与遥感图像这些技术应用中经常采用更多的级数以充分利用每个采样10或12 bits的传感器精度，并且避免计算时的近似误差。在这样的应用领域流行使用16 bits即65536个组合（或65536种颜色）。

\subsubsection{RGB颜色系统}
三原色光模式（RGB color model），又称RGB颜色模型或红绿蓝颜色模型，是一种加色模型，将红（Red）、绿（Green）、蓝（Blue）三原色的色光以不同的比例相加，以合成产生各种色彩光。

RGB颜色模型的主要目的是在电子系统中检测，表示和显示图像，比如电视和电脑，利用大脑强制视觉生理模糊化(失焦),将红绿蓝三原色子象素合成为一色彩象素,产生感知色彩(其实此真彩色并非加色法所产生的合成色彩，原因为该三原色光从来没有重叠在一起，衹是人类为了"想"看到色彩，大脑强制眼睛失焦而形成。红绿蓝三色模型在传统摄影中也有应用。在电子时代之前，基于人类对颜色的感知，RGB颜色模型已经有了坚实的理论支撑。

RGB是一种依赖于设备的颜色空间：不同设备对特定RGB值的检测和重现都不一样，因为颜色物质（荧光剂或者染料）和它们对红、绿和蓝的单独响应水平随着制造商的不同而不同，甚至是同样的设备不同的时间也不同。

\subsubsection{Gamma校正}
伽马校正（Gamma correction）又叫伽马非线性化（gamma nonlinearity）、伽马编码（gamma encoding）或是就只单纯叫伽马（gamma）。是用来针对影片或是影像系统里对于光线的辉度（luminance）或是三色刺激值（tristimulus values）所进行非线性的运算或反运算。最简单的例子中，伽马校正是由下幂定律公式所定义的：
$V_{out}=AV_{in}^\gamma$

其中A是一个常量，输入和输出的值都为非负实数值。一般地来说在A=1的通常情况下，输入输出的值的范围都是在0到1之间。伽马值$\gamma < 1$的情况有时被称作编码伽马值（encoding gamma），而执行这个编码运算所使用上述幂定律的过程也叫做伽马压缩（gamma compression）；相对地，伽马值$\gamma > 1$的情况有时也被称作解码伽马值（decoding gamma），而执行这个解码运算所使用上述幂定律的过程也叫做"伽马展开（gamma expansion）"。

\subsection{传统图像处理算法}

\subsubsection{图像增强算法}
图像增强算法用于增强对所需信息的辨别能力但同时也可能损失一些其他信息。因此在实际的应用中要对图像进行多种变换，然后再从中选取效果较好的一个。

灰度变换是指根据某种目标条件来确立变换关系然后逐点改变图像中每一个像素的灰度值。让我们假设原图像像素的灰度值D=f(x,y)，处理后图像像素的灰度值记为$D'=g(x,y)$,则灰度增强可表示为：$D'=T(D)$其中D和$D'$都在图像的灰度范围之内。函数T(D)称为灰度变换函数，它描述了输入灰度值和输出灰度值之间的转换关系。当我们确定了灰度变换函数后，一个具体的灰度增强方法便确立了。根据T(D)表达式的不同将灰度变换分为线性变换和非线性变换。

\subsubsection{图像滤波算法}
由于外部或内部的各种原因，数字图像不可避免地会有一些噪声，如拍摄图片时光照不足会产生噪点，扫描图片时我们也无法得到百分之百不含杂质的图像。这便对后续的图像处理造成了影响。因此，图像去噪是数字图像处理工作中一个必不可少的部分，对于后续的各项其他处理都有重要意义。

中值滤波是一种非线性的处理方法。常规中值滤波首先用3×3大小的滑动窗口在彩色图像上自左向右、自上向下移动，设Sxy代表以点(x,y)为中心、大小是3×3的矩形图像邻域像素值（窗口）的一组坐标。

\subsubsection{图像边缘检测算法}
图像边缘是图像的基本特征之一，它包含对人类视觉和机器视觉有价值的物体图像信息。例如在利用计算机锐化图像时便有效利用了图像的边缘信息。图像边缘检测算法是图像处理领域中一种重要的预处理技术，其广泛地应用于轮廓、特征的抽取和纹理分析等领域。

Canny边缘检测算法是针对局部区域或整幅图像确定一个门限，对图像进行分割，从而提取出边缘。Canny边缘检测算法总共有四个步骤，分别是：I.采用高斯滤波的方法对原始图像作平滑处理 II.梯度幅值与方向的计算 III.对梯度幅值的非极大值抑制IV.双阈值检测和连接边缘。

\section{次线性算法}

在大数据时代，传统的线性算法往往无法满足实时处理的需求。次线性算法的出现为解决这一问题提供了新的思路。

到这里为止，我非常感兴趣的还是那些有可能获得置信度的算法的性质。我同样确定，要寻找最优预测算法，囤积海量数据是必不可少的。

然而，并不是所有数据都有价值。就像在沙漠深处安装的监视摄像头传来的视频流那样，实际上绝大部分收集来的数据毫无意义。问题在于，在这个大数据的时代，单单读取这些没有意义的数据来验证它们，真的没有意义，所需的时间就已经不切实际了。

为了解决这个问题，理论计算机学家将注意力投向了所谓的次线性算法。这些算法的特点就是在不花时间读取所有数据的情况下，仍然能够提取数据集里最核心的内容。也就是说，这些算法能够在读取输入数据时"一目十行"。

这种算法在历史上最典型的例子就是谷歌的算法。在谷歌上搜索几乎瞬间就能得到结果，然而，谷歌的数据库包含了互联网的数百万亿网页中相当大的一部分。对于谷歌来说，在向用户返回结果之前探索整个数据库根本不在考虑范围之内，因为这样做要花上几天时间！因此，谷歌的搜索算法必须是次线性的。

谷歌用到的技巧跟图书馆和词典一样，就是预先对数据库排序和整理，迅速完成对数据库的查阅。比如说，词典将单词按照字母顺序整理，就能让使用者很快知道希望查找的单词应该会在什么地方。更厉害的是，利用字母顺序，给定词典中的一页以及要查找的单词，使用者就能知道这个单词应该在词典中这一页的前面还是后面。一般来说，在(按照全序关系)排序过的数据库中查找数据非常迅速。所谓的二分算法(大概就是你用词典查单词的时候用的方法)的计算时间是数据库大小的对数。

\subsection{傅里叶变换的优化}

在19世纪初，约瑟夫·傅里叶等人就曾研究过傅里叶变换，它是一种对声音、图像、视频以及股票走势等信号的描述方式进行变换的方法。我们在这里就只讨论声音信号。

声音可以通过鼓膜附近气压的变化来描述。但还有另一种等价的描述方式，它对于音乐的谱写来说特别简单，那就是利用组成声音的频率来描述这个声音。傅里叶变换就像一个双语词典，能够将声音的振动变化描述转换为频率描述。

实际上，傅里叶变换已经占据了我们的日常生活，据理查德·巴拉纽克教授所言，就连我们的计算机和电话每天都会进行数十亿次傅里叶变换。在每次听音乐、查看数字化图像，或者观看视频的时候，我们都在让机器进行傅里叶变换。

因此，计算傅里叶变换的算法如果有任何加速的可能性，都会在程序计算或者计算所需硬件的方面带来数十亿欧元的收益，更不用说用户节省下来的等待时间了。在1964年，詹姆斯·库利和约翰·图基就完成了一项壮举，大幅缩短了(离散)傅里叶变换的计算时间。尽管最简单的算法需要的时间与数据量的平方成正比，但库利和图基的算法，又叫快速傅里叶变换，可以在接近线性的时间内完成。也就是说，要对一百万字节(1MB)的数据进行计算的话，快速傅里叶变换只需要数百万次运算(对于现代计算机来说，这相当于几毫秒的计算)，而不是朴素算法所需的100万乘以100万次运算(差不多1分钟的计算)。

然而在大数据的时代，快速傅里叶变换还是太慢了，尤其在处理上十亿字节(GB)甚至上万亿字节(TB)的数据时更是如此。我们能不能让傅里叶变换变得比现在更快？在2012年，哈桑尼耶、因迪克、卡塔比和普赖斯发现，答案是肯定的。也可以说他们提出了一个巧妙的算法，可以对信号最主要的个频率进行近似计算，而需要的时间大概在乘以数据量的对数这个数量级。重点在于，这个叫作稀疏傅里叶变换的精妙算法的运行时间是次线性的，也就是说它几乎忽略了整个待处理的信号，却能得知这个信号大概是什么。

\section{思考的多种模式}

虽然次线性算法在计算机科学中仍不是主流，但我们的大脑似乎喜欢使用这种算法。谁又未曾一目十行读完书面材料，漫不经心地听别人讨论，或者吃饭时根本没有在意食物的味道呢？奇怪的是，有时候文章中的某些词语、讨论中的某些话题或者食物中的某些味道会吸引我们的注意力。当这种情况出现的时候，我们似乎突然就切换到了思考状态。我们会开始使用更缓慢、更精确的算法来仔细分析文章的含义、讨论的深意与本地特色菜的独特味道。

我在这里描述的，正是诺贝尔经济学奖得主丹尼尔·卡内曼的杰作《思考，快与慢》的核心。卡内曼区分了两套思考系统：系统一和系统二。系统一就像稀疏傅里叶变换，非常迅速、有效、勤奋，但有可能出现严重错误；系统二就像精确傅里叶变换，很慢且需要大量能量，懒惰但更正确。

据卡内曼所说，我们通常会将自己等同于系统二，而且几乎不会意识到系统一的存在。然而，在绝大部分时间内主导的都是系统一，而这会让我们经常犯错。试试这个：球拍和球的价格一共是1.1欧元，球拍的价格比球高1欧元，那么球的价格是多少？

很可能你立刻就看到了答案，但这个答案是错的。根据卡内曼所说，如果这种情况出现，那是因为系统一匆匆忙忙给出了第一个想到的答案；只有在之后，也许是当你看到这里的时候，你的系统二才会开始质疑系统一。

\section{迈进后严谨阶段}

数学的学习似乎也是如此进行的。在今天，如果有人问我一个关于加法、指数或者图灵机的问题，那么对我来说，不通过系统二就能回答这些问题也并非毫无可能。这是因为，长年以来我的系统二已经成功教会系统一如何在不花费脑力劳动的前提下回答这些问题！我的大脑中储存着大量捷径，能用于解决大量数学问题，即使是那些对于年轻学生来说略感困难的问题。某些人把这种能力称为数学能力，还有些人把它称为数学直觉。我更愿意把它称为系统二通过努力发现的高效捷径。

但这不是系统一的数学训练中最重要的方面。数学家、菲尔兹奖获得者陶哲轩在他的博客中将数学学习分为三个阶段，分别叫作前严谨阶段、严谨阶段和后严谨阶段。简单来说，学生首先会开始摆弄数字和数学概念，而不会忧心于自己执行的那些代数操作的有效性。然后，随着时间流逝以及接受的数学教育越来越多，严谨性的时刻就会到来，而且很快会转变为纯粹主义：一切都应该用形式化的语言做出解释。最后，后严谨阶段属于研究者，他们花上大部分时间构建不同的启发式论证，用来得出定理证明的大体轮廓，其间会暂时舍弃之前学到的形式体系。关键在于，第二个阶段似乎是到达第三个阶段的必经之路。

陶哲轩的说法可以用系统一和系统二的语言来理解。前严谨阶段对应的是系统一和系统二都没有学会严谨性的情况。当系统二发现严谨性及其重要性时，严谨阶段就开始了。在这个阶段中，系统二会教导系统一，让它意识到自身直觉的局限性。但更重要的是，系统一会由此学会测量自身的置信度，这样就能够更好地做出决定，判断自己能解决问题，还是需要向系统二求助。只有经过这种困难的学习过程，系统一才能成为系统二的完美协助者，而不会在重要的时候拖后腿。因此，成为一名优秀的数学家，基本上相当于让系统一足以在这方面胜任---但要达到这一状态，严谨阶段必不可少。

假如我能打个赌，（作为合格的贝叶斯主义者，我喜欢打赌！）我会说这就是未来人工智能的模样。

\section{实践思考}

人工智能的数学工具不仅仅是理论概念，更是解决实际问题的有力武器。在实践中，我们需要根据具体问题选择合适的数学工具：

1. **数据预处理**：使用线性代数进行数据变换和降维
2. **模型训练**：运用优化理论寻找最优参数
3. **不确定性处理**：借助概率统计量化和管理风险
4. **算法设计**：基于计算理论分析算法复杂度
5. **信息提取**：利用信息论衡量特征的重要性

线性代数在人工智能中的应用包括：
- 矩阵和向量运算：在数据表示和转换中使用，如矩阵乘法、转置和逆矩阵
- 特征值和特征向量：用于数据降维和特征选择，如主成分分析（PCA）
- 奇异值分解（SVD）：用于数据压缩、降噪和推荐系统
- 矩阵分解技术：如非负矩阵分解（NMF）用于主题建模

通过系统学习和掌握这些数学工具，我们能够更好地理解人工智能算法的本质，并在实际应用中做出明智的选择。数学不仅是人工智能的基础，更是创新的源泉。

\nocite{*}
\printbibliography[heading=bibintoc, title=\ebibname]
\newpage
